\cleardoublepage
\phantomsection
\addcontentsline{toc}{chapter}{Daftar Istilah}
\chapter*{Daftar Istilah}

\begin{longtable}{|p{3.5cm}|p{10.5cm}|}
\hline
\textbf{Istilah} & \textbf{Definisi} \\ \hline
\endfirsthead
\hline
\textbf{Istilah} & \textbf{Definisi} \\ \hline
\endhead
\hline
\endfoot
\hline
\endlastfoot
Uang & Instrumen finansial yang berfungsi sebagai penyimpan nilai dan alat tukar, dalam ekonofisika sering dimodelkan melalui analogi distribusi energi termodinamika. \\ \hline
Pasar & Tempat bertemunya penjuan dan pembeli untuk melakukan transaksi jual beli barang maupun jasa. \\ \hline
Investasi & Komitmen sejumlah modal pada aset finansial dengan ekspektasi memperoleh imbal hasil (\textit{return}) di masa depan, yang melibatkan pertukaran (\textit{trade-off}) antara risiko dan keuntungan. \\ \hline
Hamiltonian & Operator matematika yang merepresentasikan energi total sebuah sistem; dalam konteks optimasi kuantum, \textit{Hamiltonian} Ising digunakan untuk merepresentasikan fungsi biaya (\textit{cost function}). \\ \hline
Return & Tingkat pengembalian atau imbal hasil dari suatu investasi dalam periode waktu tertentu, biasanya dihitung sebagai perubahan persentase harga aset (\textit{simple return}) atau logaritma natural rasio harga (\textit{log return}). \\ \hline
Portofolio & Kumpulan atau kombinasi aset keuangan seperti saham, obligasi, dan kas yang dikelola untuk mencapai tujuan investasi tertentu melalui diversifikasi guna mereduksi risiko non-sistematis. \\ \hline
Qubit & Bit kuantum yang merupakan unit dasar informasi dalam komputasi kuantum, mampu merepresentasikan superposisi linier dari keadaan basis $|0\rangle$ dan $|1\rangle$. \\ \hline
Operator Pauli & Sekumpulan tiga matriks kompleks ($X, Y, Z$) yang bersifat \textit{Hermitian} dan \textit{unitary}, digunakan sebagai basis operator untuk memanipulasi dan mengukur keadaan \textit{qubit}. \\ \hline
Teori Prospek & Model perilaku ekonomi yang menjelaskan pengambilan keputusan dalam kondisi ketidakpastian, di mana individu cenderung lebih sensitif terhadap potensi kerugian dibandingkan keuntungan yang setara (\textit{loss aversion}). \\ \hline
Kreditur & Pihak pemberi pinjaman berupa dana atau jasa yang diharapkan akan diterima kembali pada waktu yang disepakati. \\ \hline
Debitur & Pihak yang berhutang berupa dana atau jasa kepada pihak kreditur yang wajib dilunasi sesuai kesepakatan. \\ \hline
Inflasi & Fenomena kenaikan harga barang dan jasa secara umum dan terus-menerus dalam jangka waktu tertentu yang mengakibatkan penurunan daya beli masyarakat. \\ \hline
Barren Plateau & Fenomena yang terjadi pada algoritma optimasi kuantum di mana landskap fungsi biaya menjadi sangat datar sehingga gradiennya mendekati nol, menyulitkan algoritma klasik untuk menemukan minimum global. \\ \hline
Aset & Sumber daya ekonomi yang memiliki nilai dan dapat dimiliki secara individu maupun kolektif, mencakup aset riil (real estate, komoditas) dan aset finansial (saham, obligasi). \\ \hline
Asimetri Informasi & Kondisi di mana satu pihak dalam transaksi memiliki informasi yang lebih banyak atau lebih baik daripada pihak lainnya, sehingga menciptakan ketidakseimbangan kekuatan pasar. \\ \hline
Volatilitas & Fluktuasi harga aset yang disebabkan oleh faktor-faktor sistemik yang mempengaruhi pasar secara keseluruhan, seperti kebijakan moneter, krisis finansial, atau bencana alam. \\ \hline
Likuiditas & Kemampuan suatu aset untuk diperdagangkan atau dicairkan menjadi uang tunai tanpa menyebabkan perubahan harga yang signifikan. \\ \hline
Liabilitas & Kewajiban finansial atau hutang yang harus dibayar oleh suatu entitas kepada kreditur dalam jangka waktu tertentu. \\ \hline
Herding Behavior & kecenderungan individu untuk mengikuti tindakan atau keputusan mayoritas dalam suatu kelompok, terlepas dari penilaian rasional individu tersebut terhadap informasi yang tersedia. \\ \hline

\end{longtable}
